\documentclass[11pt]{article}

\usepackage[a4paper,margin=1in]{geometry}
\usepackage[T1]{fontenc}
\usepackage{lmodern}
\usepackage{xcolor}
\usepackage{hyperref}
\usepackage{booktabs}
\usepackage{longtable}
\usepackage{tabularx}
\usepackage{enumitem}
\usepackage{listings}
\usepackage{float}
\usepackage{parskip}

\definecolor{codegray}{RGB}{245,245,245}
\definecolor{darkblue}{RGB}{25,65,105}

\hypersetup{
  colorlinks=true,
  linkcolor=darkblue,
  urlcolor=darkblue,
  pdftitle={Drone Control Center System Understanding Guide}
}

\lstset{
  basicstyle=\ttfamily\small,
  backgroundcolor=\color{codegray},
  frame=single,
  breaklines=true,
  columns=fullflexible,
  keepspaces=true,
  showstringspaces=false
}

\title{Drone Control Center\\System Understanding Guide}
\author{}
\date{\today}

\begin{document}
\maketitle
\tableofcontents
\newpage

\section{System Architecture}

This repository contains two browser applications and a shared Python laptop
backend:

\begin{itemize}[leftmargin=*]
  \item \texttt{control\_center/}: the active Drone Control Center with the
  camera, water-contact, trajectory, telemetry, control-log, and valve-trigger
  interfaces.
  \item \texttt{current\_frontend/}: the older standalone control interface
  with flight sliders, Pico outputs, telemetry, connection controls, and a
  protocol preview.
  \item \texttt{laptop\_backend/}: the WebSocket-to-UDP bridge between the
  browser and Raspberry Pi.
\end{itemize}

\subsection{Runtime data flow}

\begin{lstlisting}
Synthetic or real camera producer
        |
        | WebSocket JSON, normally ws://127.0.0.1:8765
        v
Laptop backend: laptop_backend/control_center_backend.py
        ^
        | WebSocket JSON, ws://127.0.0.1:5173
        v
Browser UI: control_center/
        |
        | UDP JSON controls and triggers, port 5000
        v
Raspberry Pi gateway
        |                    |
        | MAVLink             | SPI
        v                    v
Flight controller          Pico
        ^
        |
        | UDP JSON telemetry, port 5001
        |
Laptop backend
\end{lstlisting}

The browser does not communicate directly with the Raspberry Pi. It connects
to the laptop backend over WebSocket. The backend validates and forwards
control and trigger messages as UDP JSON packets, and broadcasts incoming Pi
telemetry back to connected browsers.

\subsection{Video path}

The video producer sends a \texttt{VIDEO\_FRAME} message to the backend:

\begin{lstlisting}
{
  "type": "VIDEO_FRAME",
  "seq": 12,
  "time_ms": 1727483000123,
  "content_type": "image/jpeg",
  "data": "<base64 JPEG>",
  "water_contact": true,
  "sensors": {
    "water_detected": true
  }
}
\end{lstlisting}

The backend validates JPEG, PNG, or WebP frames and relays them to browser
clients through port 5173. The active frontend converts the base64 payload to
a browser data URL and displays it. The water-contact box becomes green when
either \texttt{water\_contact} or \texttt{sensors.water\_detected} is true.

The bundled Gaussian-noise producer is only a development/test camera. It
defaults to reporting water detected and supports
\texttt{--no-water-detected} for testing the inactive state.

\subsection{Control path}

The browser sends \texttt{CONTROL} messages containing normalized flight axes
and four Pico outputs:

\begin{lstlisting}
{
  "version": 1,
  "type": "CONTROL",
  "flight": {
    "roll": 0.2,
    "pitch": -0.1,
    "yaw": 0.05,
    "throttle": 0.3
  },
  "pico": {
    "output0": 100,
    "output1": 0,
    "output2": -200,
    "output3": 50
  }
}
\end{lstlisting}

The backend validates flight values from \(-1.0\) to \(1.0\), requires
non-negative throttle, validates Pico values from \(-1000\) to \(1000\), adds
its own sequence number, and forwards the packet to the Pi.

The active keyboard mapping is:

\begin{center}
\begin{tabular}{ll}
\toprule
\textbf{Input} & \textbf{Control}\\
\midrule
\texttt{W} & Positive pitch\\
\texttt{S} & Negative pitch\\
\texttt{A} & Negative roll\\
\texttt{D} & Positive roll\\
Up arrow & Positive throttle\\
Down arrow & Negative throttle, subject to validation\\
\bottomrule
\end{tabular}
\end{center}

The browser emits controls at 20~Hz. This is a human-in-the-loop interface,
not a hard-real-time flight stabilization loop.

\subsection{Valve trigger path}

The active interface has a 2-by-2 valve grid:

\begin{center}
\begin{tabular}{ll}
\toprule
Valve 1 ON & Valve 2 ON\\
Valve 1 OFF & Valve 2 OFF\\
\bottomrule
\end{tabular}
\end{center}

The buttons send \texttt{TRIGGER} messages with one of:
\texttt{valve\_1\_on}, \texttt{valve\_2\_on},
\texttt{valve\_1\_off}, or \texttt{valve\_2\_off}. The Pi gateway must map
these names to the appropriate GPIO, Pico, or actuator operation.

\subsection{Telemetry path}

The Pi sends \texttt{TELEMETRY} JSON datagrams to the laptop's UDP port 5001.
The backend adds \texttt{bridge\_received\_time\_ms}, tracks telemetry age,
and broadcasts the message to browser clients.

Telemetry may contain flight mode, armed state, attitude, acceleration,
angular velocity, position, battery, ultrasonic readings, water status, and
other sensors. The active UI uses the flight-controller fields for the IMU,
trajectory, and status panels while preserving additional fields in transit.

\section{File Responsibilities}

\subsection{\texttt{laptop\_backend/}}

\begin{longtable}{p{0.32\textwidth}p{0.62\textwidth}}
\toprule
\textbf{File} & \textbf{Responsibility}\\
\midrule
\texttt{control\_center\_backend.py} &
Main laptop bridge. Runs the browser WebSocket server, validates controls and
triggers, sends UDP commands to the Pi, receives UDP telemetry, connects to
the video producer, validates frames, relays video, and reports bridge status.\\
\texttt{test\_control\_center\_io.py} &
Bench-test harness. Sends neutral controls and synthetic telemetry, receives
backend messages, and records JSON Lines logs. It does not control real
hardware.\\
\texttt{requirements.txt} &
Python dependency declaration, primarily the \texttt{websockets} package.\\
\texttt{.env.example} &
Example network settings and configuration reference.\\
\texttt{control\_center\_io.jsonl} &
Generated test log, not application source.\\
\bottomrule
\end{longtable}

Important backend options are:

\begin{lstlisting}
--pi-ip             Raspberry Pi address; default 192.168.1.50
--pi-control-port   Pi command UDP port; default 5000
--telemetry-port   Laptop telemetry UDP port; default 5001
--host              Browser WebSocket bind host; default 127.0.0.1
--port              Browser WebSocket port; default 5173
--video-url         Camera producer WebSocket; default ws://127.0.0.1:8765
\end{lstlisting}

\subsection{\texttt{current\_frontend/}}

This is the older React/Vite application.

\begin{longtable}{p{0.32\textwidth}p{0.62\textwidth}}
\toprule
\textbf{File} & \textbf{Responsibility}\\
\midrule
\texttt{src/main.jsx} & React bootstrap and root mounting.\\
\texttt{src/App.jsx} &
Connection form, automatic reconnect, flight sliders, Pico outputs,
telemetry, link statistics, protocol preview, and event log.\\
\texttt{src/styles.css} & Layout, cards, controls, telemetry, and status styles.\\
\texttt{index.html} & Vite HTML entry point.\\
\texttt{vite.config.js} & Binds this application to \texttt{127.0.0.1:6970}.\\
\texttt{package.json} & npm scripts and React/Vite dependencies.\\
\bottomrule
\end{longtable}

Its default backend URL is \texttt{ws://127.0.0.1:5173}, stored in browser
local storage under \texttt{wsUrl}. This frontend is not the active
camera/water/valve interface.

\subsection{\texttt{control\_center/}}

\begin{longtable}{p{0.32\textwidth}p{0.62\textwidth}}
\toprule
\textbf{File} & \textbf{Responsibility}\\
\midrule
\texttt{src/main.tsx} &
Active UI, backend WebSocket client, video rendering, water signal handling,
IMU and trajectory views, keyboard controls, telemetry, logs, and valve
triggers.\\
\texttt{src/style.css} &
Full-screen two-column layout, camera, IMU, trajectory, water panel, and
2-by-2 valve grid styling.\\
\texttt{gaussian\_video\_server.py} &
Synthetic JPEG video WebSocket producer with optional water detection signal.\\
\texttt{index.html} & Vite HTML entry point.\\
\texttt{vite.config.ts} & Binds the active frontend to \texttt{127.0.0.1:6969}.\\
\texttt{package.json} & npm scripts and active React/Vite dependencies.\\
\texttt{requirements.txt} & Python dependencies for the synthetic video server.\\
\bottomrule
\end{longtable}

\section{Ports and Protocols}

\begin{longtable}{rllp{0.43\textwidth}}
\toprule
\textbf{Port} & \textbf{Transport} & \textbf{Endpoints} & \textbf{Meaning}\\
\midrule
6969 & TCP/HTTP & Vite $\rightarrow$ browser &
Serves the active \texttt{control\_center/} application.\\
6970 & TCP/HTTP & Vite $\rightarrow$ browser &
Serves the older \texttt{current\_frontend/} application.\\
5173 & TCP/WebSocket & Browser $\leftrightarrow$ laptop &
Controls, triggers, bridge status, telemetry, errors, and relayed video.\\
8765 & TCP/WebSocket & Producer $\rightarrow$ laptop &
Synthetic or real camera \texttt{VIDEO\_FRAME} source.\\
5000 & UDP & Laptop $\rightarrow$ Pi &
Validated \texttt{CONTROL} and \texttt{TRIGGER} JSON datagrams.\\
5001 & UDP & Pi $\rightarrow$ laptop &
\texttt{TELEMETRY} JSON datagrams.\\
\bottomrule
\end{longtable}

Port 6969 and 6970 are HTTP development servers, not control ports. Port
5173 is the browser protocol endpoint. Port 8765 is used by the backend as a
camera input. UDP is connectionless, so it does not guarantee delivery,
ordering, or retransmission. Sequence numbers, timestamps, and Pi-side
watchdogs are therefore important.

Individual WebSocket messages must remain below the backend's 2~MiB message
limit. Base64 video increases payload size, so real camera resolutions and
quality settings may need to be reduced.

\subsection{Port and address reference}

The active \texttt{control\_center/} frontend is served at
\texttt{http://127.0.0.1:6969}. The older \texttt{current\_frontend/}
frontend is served at \texttt{http://127.0.0.1:6970}. In the current
configuration, both browser applications connect to the laptop backend at
\texttt{ws://127.0.0.1:5173}; these loopback addresses refer to services on
the laptop itself.

The complete data flow is:

\begin{lstlisting}
Browser
  | HTTP :6969 or :6970
  v
Frontend application
  | WebSocket ws://127.0.0.1:5173
  v
Laptop backend
  | UDP 192.168.1.50:5000
  v
Raspberry Pi

Raspberry Pi
  | UDP <laptop-Ethernet-IP>:5001
  v
Laptop backend
  | WebSocket broadcast :5173
  v
Browser

Video producer
  | WebSocket :8765
  v
Laptop backend
\end{lstlisting}

For a direct Ethernet connection, a typical static configuration is
\texttt{192.168.1.10} for the laptop Ethernet adapter and
\texttt{192.168.1.50} for the Raspberry Pi, with subnet mask
\texttt{255.255.255.0}. The backend sends control and trigger packets to the
Pi at \texttt{192.168.1.50:5000}. The Pi sends telemetry to the laptop's
Ethernet address, for example \texttt{192.168.1.10:5001}; it must not send
telemetry to \texttt{127.0.0.1}, which refers to the Pi itself.

The backend's default video URL is \texttt{ws://127.0.0.1:8765}, which is
appropriate when the producer runs on the laptop. If the producer runs on the
Pi, configure \texttt{--video-url ws://192.168.1.50:8765} instead. The
laptop firewall must allow inbound UDP traffic on port 5001 for Pi telemetry.

\section{Test Scripts and Functions}

\subsection{\texttt{test\_control\_center\_io.py}}

Run it from \texttt{laptop\_backend/} after starting the backend:

\begin{lstlisting}
..\.venv\Scripts\python.exe test_control_center_io.py ^
  --duration 10 ^
  --rate-hz 10 ^
  --output .\control_center_io.jsonl
\end{lstlisting}

\begin{description}[leftmargin=!,labelwidth=0.28\textwidth]
\item[\texttt{EventLog}] Opens the JSONL file, writes timestamped events, and
closes the file.
\item[\texttt{telemetry\_packet}] Creates deterministic changing telemetry
using sine and cosine values for attitude, position, and sensors.
\item[\texttt{control\_packet}] Creates neutral flight and Pico controls.
\item[\texttt{send\_telemetry}] Sends synthetic telemetry UDP datagrams at
the selected rate and duration.
\item[\texttt{exercise\_websocket}] Connects to the backend, sends
\texttt{CONFIG} and controls, receives messages, and logs all activity.
\item[\texttt{run}] Runs the WebSocket and UDP tests concurrently.
\end{description}

This harness tests control, telemetry, and logging. It does not test real
camera hardware or physical valve actuation.

\subsection{\texttt{gaussian\_video\_server.py}}

\begin{lstlisting}
..\.venv\Scripts\python.exe gaussian_video_server.py
..\.venv\Scripts\python.exe gaussian_video_server.py --no-water-detected
\end{lstlisting}

The main functions are:

\begin{description}[leftmargin=!,labelwidth=0.28\textwidth]
\item[\texttt{create\_noise\_frame}] Generates and base64-encodes one JPEG.
\item[\texttt{stream\_video}] Sends frames at the configured FPS.
\item[\texttt{handle\_client}] Handles one connected backend client.
\item[\texttt{run\_server}] Starts the video WebSocket server.
\item[\texttt{parse\_args}] Parses host, port, dimensions, FPS, and water state.
\end{description}

This is a development producer, not a real camera driver.

\section{Complete Run Order}

\subsection{First-time setup}

\begin{lstlisting}
cd C:\Users\ASUS\Desktop\CMU\Sem1\EMSD\Frontend\drone-control-center
python -m venv .venv
.\.venv\Scripts\Activate.ps1
python -m pip install -r laptop_backend\requirements.txt
python -m pip install -r control_center\requirements.txt
\end{lstlisting}

Install JavaScript dependencies in both applications:

\begin{lstlisting}
cd control_center
npm install
cd ..\current_frontend
npm install
\end{lstlisting}

\subsection{Start services}

\begin{enumerate}
\item Start the synthetic video producer:
\begin{lstlisting}
cd control_center
..\.venv\Scripts\python.exe gaussian_video_server.py
\end{lstlisting}

\item Start the laptop backend:
\begin{lstlisting}
cd ..\laptop_backend
..\.venv\Scripts\python.exe control_center_backend.py --pi-ip 127.0.0.1
\end{lstlisting}
For real hardware, replace \texttt{127.0.0.1} with the Pi address.

\item Start the active UI:
\begin{lstlisting}
cd ..\control_center
npm run dev
\end{lstlisting}
Open \url{http://127.0.0.1:6969}.

\item Optionally start the older UI:
\begin{lstlisting}
cd ..\current_frontend
npm run dev
\end{lstlisting}
Open \url{http://127.0.0.1:6970} and verify its WebSocket URL is
\texttt{ws://127.0.0.1:5173}.

\item Optionally run the I/O harness:
\begin{lstlisting}
cd ..\laptop_backend
..\.venv\Scripts\python.exe test_control_center_io.py --duration 10 --rate-hz 10
\end{lstlisting}
\end{enumerate}

\subsection{Safe bench-test order}

\begin{enumerate}
\item Remove propellers and keep the vehicle disarmed.
\item Start the Pi gateway with actuators disabled.
\item Start the backend with the correct Pi address.
\item Verify browser-to-backend connectivity.
\item Verify telemetry and increasing sequence numbers.
\item Test safe Pico outputs.
\item Test valves in a physically safe setup.
\item Connect flight-controller functionality while disarmed.
\item Add real video and validate the camera path.
\end{enumerate}

\section{Interfacing with a Raspberry Pi}

\subsection{Pi gateway responsibilities}

The Raspberry Pi needs a gateway program that:

\begin{enumerate}
\item Listens for UDP JSON commands on port 5000.
\item Parses and validates \texttt{CONTROL} and \texttt{TRIGGER} packets.
\item Maps flight commands to MAVLink or the selected flight API.
\item Maps Pico values to SPI.
\item Maps valve names to GPIO, Pico, or actuator operations.
\item Reads flight-controller, GPS, battery, ultrasonic, and water sensors.
\item Sends \texttt{TELEMETRY} JSON to the laptop on UDP port 5001.
\item Enforces a control timeout and safe actuator state.
\end{enumerate}

The Pi should own safety-critical behavior. The browser and laptop are
operator and networking layers, not the stabilization loop.

\subsection{Network configuration}

Give the Pi a reachable address, preferably with a DHCP reservation or static
configuration. Start the backend with:

\begin{lstlisting}
..\.venv\Scripts\python.exe control_center_backend.py ^
  --pi-ip 192.168.1.50 ^
  --pi-control-port 5000 ^
  --telemetry-port 5001
\end{lstlisting}

The laptop firewall must allow inbound UDP port 5001. The Pi must be able to
send telemetry to the laptop's IP address.

\subsection{Pi command receiver}

The Pi gateway should dispatch by message type:

\begin{lstlisting}[language=Python]
if message["type"] == "CONTROL":
    # Validate version, ranges, sequence, timestamp, and freshness.
    # Send flight values to the flight controller.
    # Send Pico values over SPI.
elif message["type"] == "TRIGGER":
    # Dispatch valve_1_on, valve_2_on,
    # valve_1_off, and valve_2_off.
\end{lstlisting}

Recommended Pi-side checks include:

\begin{itemize}
\item Reject unknown protocol versions.
\item Reject non-finite, out-of-range, or stale values.
\item Track sequence numbers and detect packet gaps.
\item Use a watchdog timeout and return actuators to a safe state.
\item Keep emergency behavior local to the Pi.
\end{itemize}

\subsection{Pi telemetry}

The Pi should send packets containing at least:

\begin{lstlisting}
{
  "version": 1,
  "type": "TELEMETRY",
  "seq": 42,
  "time_ms": 1727483000123,
  "fc": {
    "flight_mode": "STABILIZE",
    "armed": false,
    "attitude": {"roll": 0.0, "pitch": 0.0, "yaw": 0.0},
    "position": {"lat": 40.4432, "lon": -79.9436, "alt": 12.0},
    "battery": {"voltage_v": 15.8, "remaining_percent": 87}
  },
  "water_contact": false,
  "sensors": {"water_detected": false}
}
\end{lstlisting}

Send this datagram to \texttt{<laptop-ip>:5001}. The backend forwards it to
all browser clients.

\subsection{Real camera integration}

A real camera producer should capture frames, encode them as JPEG, PNG, or
WebP, base64-encode the image, and publish \texttt{VIDEO\_FRAME} messages to a
WebSocket endpoint. Configure the backend with:

\begin{lstlisting}
..\.venv\Scripts\python.exe control_center_backend.py ^
  --pi-ip 192.168.1.50 ^
  --video-url ws://192.168.1.50:8765
\end{lstlisting}

Keep serialized messages below 2~MiB. If full-resolution frames exceed this
limit, reduce resolution or quality, or replace base64 JSON with a video
transport designed for continuous media.

\subsection{Recommended integration sequence}

\begin{enumerate}
\item Verify the laptop backend and browser locally.
\item Verify the I/O harness and synthetic telemetry.
\item Implement a Pi UDP receiver and send harmless telemetry.
\item Confirm telemetry appears in the UI.
\item Send neutral controls and verify Pi sequence numbers.
\item Test safe Pico output values.
\item Implement and test valve mapping.
\item Connect the flight controller while disarmed.
\item Add real camera streaming.
\item Validate watchdog and failsafe behavior.
\end{enumerate}

\end{document}
